\subsection{Постановка целей и задач по построению системы редактирования видео по тексту}

Следовательно, крайне важно создать системы, которые могут эффективно использовать уже существующие модели нейронных сетей в применении задачи редактирования видео. Давайте исследуем потенциальное решение для повышения скорости моделей нейронных сетей путем разработки оптималной архитерктуры и оптимизации моделей.

Цель работы – разработка методов динамической замены и модификации объектов в видеоданных во временных рамках приближенных к реальному времени

Для достижения поставленной цели в работе были решены следующие задачи:
\begin{itemize}
    \item Выбрать оптимальную базовую архитектуру для решения данной задачи.
    \item Оптимизация выбранной архитектуры для ускорения работы модели.
    \item Верификация решения и полученные метрики.
\end{itemize}

Решив данные задачи как результат работы будет получен: нейросетевая моделей для редактирования видео по текстовому запросу.
